早泄的诊断以\textbf{病史采集为核心}，量表与检查用于量化与鉴别，而非"一次测定终身"。

\subsubsection{早泄诊断工具（PEDT）}

PEDT（Premature Ejaculation Diagnostic Tool）是国际通用的 5 条目自评量表，已有中文版本，广泛应用于门诊与科研。5 个条目分别评估：射精控制困难的程度、射精是否早于期望、是否受轻微刺激即射精、是否因射精过早而沮丧、是否担心射精过早影响伴侣满意度；每题按 0--4 分计，总分为 0--20 分。

\begin{itemize}
	\item 总分 \textbf{≤8 分}：不支持早泄
	\item \textbf{9--10 分}：疑似早泄，需结合临床判断
	\item \textbf{≥11 分}：支持早泄的诊断
\end{itemize}

\subsubsection{阴道内射精潜伏期（IELT）的测量}

IELT（Intravaginal Ejaculatory Latency Time）指阴茎插入阴道到射精所经历的时间，是量化早泄的核心指标。规范测量采用\textbf{秒表法}，由伴侣计时，连续记录多次（一般不少于 4 次）取中位数，以减小单次波动的影响。凭记忆自我估计的时间往往偏短，不能代替客观记录，但也不必强求分秒精确。

\subsubsection{体格检查与实验室检查}

\begin{itemize}
	\item \textbf{生殖系统检查}：评估包皮状态（包茎或过长）、龟头敏感度、阴茎硬结，必要时做前列腺触诊
	\item \textbf{基础实验室检查}：必要时检测甲状腺功能（甲亢可致早泄）、性激素（睾酮）、血糖及前列腺相关指标
	\item \textbf{合并问题筛查}：评估是否合并勃起功能障碍、慢性前列腺炎、焦虑或抑郁——ED 与 PE 常同时存在，须一并处理
\end{itemize}

\begin{tcolorbox}[colback=pink!5!white,colframe=red!50!black,title=贴心小叮咛]
量表是"体检表"而不是"考试卷"。不少人只因看过夸张的影视内容，就自认"早泄"，这与真正的 PE 相去甚远。评估的意义，正是把"对时间的焦虑"与"真正的射精控制障碍"区分开来。
\end{tcolorbox}
